\documentclass[9pt]{article}

% --- Packing + layout ---
\usepackage{fancyhdr}
% Header setup moved after \begin{document}
\usepackage[letterpaper,margin=0.3in,includehead,includefoot,headheight=14pt,headsep=6pt]{geometry} % very small margins + header/footer space
\setlength{\columnsep}{0.15in}                  % tighter gap between columns
\usepackage{microtype}                           % better glyph packing
\usepackage{multicol}
%\pagestyle{empty}                                % no header/footer/page numbers
\setlength{\parindent}{0pt}
\setlength{\parskip}{0pt}

% Tighten section spacing
\usepackage[compact,small]{titlesec}
\titlespacing*{\section}{0pt}{2pt}{1pt}
\titlespacing*{\subsection}{0pt}{2pt}{1pt}

% Tighten lists
\usepackage{enumitem}
\setlist{nosep,leftmargin=*}
\setlist[itemize]{noitemsep,topsep=0pt,parsep=0pt,partopsep=0pt,leftmargin=*}
\setlist[enumerate]{noitemsep,topsep=0pt,parsep=0pt,partopsep=0pt,leftmargin=*}
% ---- Problem labeling helpers ----
\newcommand{\exliststart}[1]{\begin{enumerate}[label=\textbf{(\Alph*)},leftmargin=*,itemsep=2pt,topsep=2pt]\def\currentex{#1}}
\newcommand{\exitem}{\item \textit{[\currentex]} }

% Tighten math spacing
\setlength{\abovedisplayskip}{2pt}
\setlength{\belowdisplayskip}{2pt}
\setlength{\abovedisplayshortskip}{1pt}
\setlength{\belowdisplayshortskip}{1pt}

% Math + links
\usepackage{amsmath,amssymb}
\usepackage[hidelinks]{hyperref}

% Callout boxes
\usepackage[most]{tcolorbox}
\tcbset{colback=white,colframe=black,boxrule=0.4pt,arc=2pt,left=6pt,right=6pt,top=4pt,bottom=4pt}
\newtcolorbox{notebox}{title=Note}
\newtcolorbox{solutionbox}{title=Solution}

% --- Fancy header/footer (moved to preamble and applied to 'plain' pages too) ---
\setlength{\headheight}{14pt}
\pagestyle{fancy}
\fancyhf{}
\fancyhead[L]{\textbf{Core Assessment 1 2026-01-05}}
\fancyfoot[C]{\thepage}
\renewcommand{\headrulewidth}{0.4pt}
\fancypagestyle{plain}{%
  \fancyhf{}
  \fancyhead[L]{\textbf{Core Assessment 1}}
  \fancyfoot[C]{\thepage}
  \renewcommand{\headrulewidth}{0.4pt}
}


\begin{document}

The following problems from Rosen are in a pool of questions that will determine the makeup of your C1 Assessment!

\textbf{Please Practice these problems!}

\begin{multicols}{3}

\section*{Construct truth tables and verify logical identities (1.1, 1.3)}

\subsection*{A.1 (1.1 Exercise 31)}
Construct a truth table for each of these compound propositions.
\begin{itemize}[noitemsep,topsep=0pt,parsep=0pt,partopsep=0pt]
  \item $p \lor \neg p$
  \item $(p \lor \neg q) \to q$
\end{itemize}

\subsection*{A.2 (1.1 Exercise 31)}
Construct a truth table for each of these compound propositions.
\begin{itemize}[noitemsep,topsep=0pt,parsep=0pt,partopsep=0pt]
  \item $(p \to q) \leftrightarrow (\neg q \to \neg p)$
  \item $(p \to q) \to (q \to p)$
\end{itemize}

\subsection*{A.3 (1.1 Exercise 35)}
Construct a truth table for each of these compound propositions.
\begin{itemize}[noitemsep,topsep=0pt,parsep=0pt,partopsep=0pt]
  \item $\neg p \leftrightarrow q$
  \item $(p \to q) \lor (\neg p \to q)$
\end{itemize}

\subsection*{A.4 (1.1 Exercise 35)}
Construct a truth table for each of these compound propositions.
\begin{itemize}[noitemsep,topsep=0pt,parsep=0pt,partopsep=0pt]
  \item $(p \leftrightarrow q) \lor (\neg p \leftrightarrow q)$
  \item $(\neg p \leftrightarrow \neg q) \leftrightarrow (p \leftrightarrow q)$
\end{itemize}

\subsection*{A.5 (1.1 Exercise 37)}
Construct a truth table for each of these compound propositions.
\begin{itemize}[noitemsep,topsep=0pt,parsep=0pt,partopsep=0pt]
  \item $\neg p \to (q \to r)$
  \item $(p \to q) \lor (\neg p \to r)$
\end{itemize}

\subsection*{A.6 (1.1 Exercise 37)}
Construct a truth table for each of these compound propositions.
\begin{itemize}[noitemsep,topsep=0pt,parsep=0pt,partopsep=0pt]
  \item $(p \leftrightarrow q) \lor (\neg q \leftrightarrow r)$
  \item $(\neg p \leftrightarrow \neg q) \leftrightarrow (q \leftrightarrow r)$
\end{itemize}


\section*{Translate between natural language statements and propositional/predicate logic (1.1, 1.4)}

\subsection*{B.1 (1.4 Exercise 9)}
Let $P(x)$ be the statement "x can speak Russian" and let $Q(x)$ be the statement "x knows the computer language C++." Express each of these sentences in terms of $P(x)$, $Q(x)$, quantifiers, and logical connectives. The domain for quantifiers consists of all students at your school.
\begin{itemize}[noitemsep,topsep=0pt,parsep=0pt,partopsep=0pt]
  \item a) There is a student at your school who can speak Russian and who knows C++.
  \item b) There is a student at your school who can speak Russian but who doesn’t know C++.
  \item c) Every student at your school either can speak Russian or knows C++.
  \item d) No student at your school can speak Russian or knows C++.
\end{itemize}

\subsection*{B.2 (1.4 Exercise 10)}
Let $C(x)$ be the statement "x has a cat," let $D(x)$ be the statement "x has a dog," and let $F(x)$ be the statement "x has a ferret." Express each of these statements in terms of $C(x)$, $D(x)$, $F(x)$, quantifiers, and logical connectives. Let the domain consist of all students in your class.
\begin{itemize}[noitemsep,topsep=0pt,parsep=0pt,partopsep=0pt]
  \item a) A student in your class has a cat, a dog, and a ferret.
  \item b) All students in your class have a cat, a dog, or a ferret.
  \item c) Some student in your class has a cat and a ferret, but not a dog.
\end{itemize}
\subsection*{B.3 (1.4 Exercise 23)}
Translate in two ways each of these statements into logical expressions using predicates, quantifiers, and logical connectives. First, let the domain consist of the students in your class and second, let it consist of all people.
\begin{itemize}[noitemsep,topsep=0pt,parsep=0pt,partopsep=0pt]
  \item a) There is a person in your class who was not born in California.
  \item b) A student in your class has been in a movie.
  \item c) No student in your class has taken a course in logic programming.
\end{itemize}

\subsection*{B.4 (1.4 Exercise 24)}
Translate in two ways each of these statements into logical expressions using predicates, quantifiers, and logical connectives. First, let the domain consist of the students in your class and second, let it consist of all people.
\begin{itemize}[noitemsep,topsep=0pt,parsep=0pt,partopsep=0pt]
  \item a) Everyone in your class has a cellular phone.
  \item b) Somebody in your class has seen a foreign movie.
  \item c) There is a person in your class who cannot swim.
\end{itemize}

\subsection*{B.5 (1.4 Exercise 30)}
Suppose the domain of the propositional function $P(x, y)$ consists of pairs $x$ and $y$, where $x$ is 1, 2, or 3 and $y$ is 1, 2, or 3. Write out these propositions using disjunctions and conjunctions.
\begin{itemize}[noitemsep,topsep=0pt,parsep=0pt,partopsep=0pt]
  \item a) $\exists x\, P(x, 3)$ 
  \item b) $\forall y\, P(1, y)$
  \item c) $\exists y\, \neg P(2, y)$ 
  \item d) $\forall x\, \neg P(x, 2)$
\end{itemize}

\subsection*{B.6 (1.4 Exercise 59)}
Let $P(x)=$"x is a professor", $I(x)=$ "x is ignorant", and $V(x)=$"x is vain". Express each of these statements using quantifiers; logical connectives; and $P(x)$, $I(x)$, and $V(x)$, where the domain consists of all people.
\begin{itemize}[noitemsep,topsep=0pt,parsep=0pt,partopsep=0pt]
  \item a) No professors are ignorant.
  \item b) All ignorant people are vain.
  \item c) No professors are vain.
\end{itemize}

\subsection*{B.7 (1.4 Exercise 60)}
Let $C(x)=$"x is a clear explanation", $S(x)=$"x is satisfactory,", and $E(x)=$"x is an excuse,". Suppose that the domain for $x$ consists of all English text. Express each of these statements using quantifiers, logical connectives, and $C(x)$, $S(x)$, and $E(x)$. The domain for quantifiers consists of all things.
\begin{itemize}[noitemsep,topsep=0pt,parsep=0pt,partopsep=0pt]
  \item a) All clear explanations are satisfactory.
  \item b) Some excuses are unsatisfactory.
  \item c) Some excuses are not clear explanations.
\end{itemize}

\section*{Negate quantified statements (1.4)}

\subsection*{C.1 (1.4 Exercise 32)}
Express each of these statements using quantifiers. Then form the negation of the statement so that no negation is to the left of a quantifier. The domain for quantifiers consists of all things.
\begin{itemize}[noitemsep,topsep=0pt,parsep=0pt,partopsep=0pt]
  \item \textbf{a)} All dogs have fleas.
  \item \textbf{b)} There is a horse that can add.
  \item \textbf{c)} Every koala can climb.
  \item \textbf{d)} No monkey can speak French.
  \item \textbf{e)} There exists a pig that can swim and catch fish.
\end{itemize}

\subsection*{C.2 (1.4 Exercise 33)}
Express each of these statements using quantifiers. Then form the negation of the statement, so that no negation is to the left of a quantifier. The domain for quantifiers consists of all things.
\begin{itemize}[noitemsep,topsep=0pt,parsep=0pt,partopsep=0pt]
  \item \textbf{a)} Some old dogs can learn new tricks.
  \item \textbf{b)} No rabbit knows calculus.
  \item \textbf{c)} Every bird can fly.
  \item \textbf{d)} There is no dog that can talk.
  \item \textbf{e)} There is no one in this class who knows French and Russian.
\end{itemize}

\subsection*{C.3 (1.4 Exercise 34)}
Express the negation of each of these statements using quantifiers. Then form the negation of the statement, so that no negation is to the left of a quantifier. The domain for quantifiers consists of all things.

\begin{itemize}[noitemsep,topsep=0pt,parsep=0pt,partopsep=0pt]
  \item \textbf{a)} Some drivers do not obey the speed limit.
  \item \textbf{b)} All Swedish movies are serious.
  \item \textbf{c)} No one can keep a secret.
\end{itemize}

\subsection*{C.4 (1.4 Exercise 35)}
Find a counterexample, if possible, to these universally quantified statements, where the domain for all variables consists of all integers.
\begin{itemize}[noitemsep,topsep=0pt,parsep=0pt,partopsep=0pt]
  \item \textbf{a)} $\forall x\, (x^2 \ge x)$
  \item \textbf{b)} $\forall x\, (x > 0 \lor x < 0)$
  \item \textbf{c)} $\forall x\, (x = 1)$
\end{itemize}

\section*{Recognize tautologies, contradictions, contingencies (1.2/1.3)}

\subsection*{D.1 (1.2 Exercise 9)}
Determine if the system is consistent or not. If it is consistent (satisfiable), state an assignment that satisfies all conditions.
\newline
"The system is in multiuser state if and only if it is operating normally. If the system is operating normally, the kernel is functioning. The kernel is not functioning or the system is in interrupt mode. If the system is not in multiuser state, then it is in interrupt mode. The system is not in interrupt mode."

\subsection*{D.2 (1.2 Exercise 10)}
Determine if the system is consistent (satisfiable) or not. If it is consistent, state an assignment that satisfies all conditions.
\begin{itemize}[noitemsep,topsep=0pt,parsep=0pt,partopsep=0pt]
  \item "Whenever the system software is being upgraded, users cannot access the file system."
  \item "If users can access the file system, then they can save new files."
  \item "If users cannot save new files, then the system software is not being upgraded."
\end{itemize}

\subsection*{D.3 (1.2 Exercise 11)}
Determine if the system is consistent (satisfiable) or not. If it is consistent, state an assignment that satisfies all conditions.
\newline
"The router can send packets to the edge system only if it supports the new address space. For the router to support the new address space, it is necessary that the latest software release be installed. The router does not support the new address space."

\subsection*{D.4 (1.2 Exercise 12)}
Determine if the system is consistent (satisfiable) or not. If it is consistent, state an assignment that satisfies all conditions.
\begin{itemize}[noitemsep,topsep=0pt,parsep=0pt,partopsep=0pt]
  \item "If the file system is not locked, then new messages will be queued."
  \item "If the file system is not locked, then the system is functioning normally, and conversely."
  \item "If new messages are not queued, then they will be sent to the message buffer."
  \item "If the file system is not locked, then new messages will be sent to the message buffer."
  \item "New messages will not be sent to the message buffer."
\end{itemize}

\section*{Work with conditionals, biconditionals, and their equivalences (1.1, 1.3)}


\subsection*{E.3 (1.3 Exercise 22)}
Show that $(p \rightarrow q) \land (p \rightarrow r)$ and $p \rightarrow (q \land r)$ are logically equivalent using logical equivalences.

\subsection*{E.4 (1.3 Exercise 23)}
Show that $(p \rightarrow r) \land (q \rightarrow r)$ and $(p \lor q) \rightarrow r$ are logically equivalent using logical equivalences.

\subsection*{E.5 (1.3 Exercise 24)}
Show that $(p \rightarrow q) \lor (p \rightarrow r)$ and $p \rightarrow (q \lor r)$ are logically equivalent using logical equivalences.

\end{multicols}
\end{document}
